% Oberwolfach Arbeitsgemeinschaft 2642 -- abstract and introduction by the organizers
% Draft. Compiles standalone with amsart; paste the abstract / MSC / introduction
% into the OWR report template. References are in owr_introduction.bib.

\documentclass{amsart}
\usepackage{amssymb}
\usepackage[T1]{fontenc}

\newcommand{\PP}{\mathbb{P}}
\newcommand{\cO}{\mathcal{O}}
\newcommand{\cF}{\mathcal{F}}
\newcommand{\cX}{\mathcal{X}}
\newcommand{\Db}{D^{b}}
\newcommand{\DCox}{D_{\mathrm{Cox}}}
\DeclareMathOperator{\Cl}{Cl}

\begin{document}

\title{Arbeitsgemeinschaft: New Homological Methods in Toric Geometry}
\author{Michael K.\ Brown, David Eisenbud, Daniel Erman}
\date{11--16 October 2026}
\maketitle

%--------------------------------------------------------------------
\section*{Abstract}

This Arbeitsgemeinschaft gave an account of recent advances in the
homological study of toric varieties. The main themes of the lectures were (1) the Hanlon--Hicks--Lazarev
resolution of toric subvarieties by line bundles and its applications to conjectures of
Bondal, Berkesch--Erman--Smith, and Orlov; (2) the Cox
category, which glues together the derived categories of all toric stacks sharing a Cox
ring and in which the Bondal--Thomsen collection is a full strong exceptional collection
(a revised form of King's conjecture); and (3) Tate resolutions over toric varieties.

\medskip
\noindent\textbf{Mathematics Subject Classification (2020):}
14M25, 14F08, 13D02, 18G80, 13D09, 16S37.

%--------------------------------------------------------------------
\section*{Introduction by the Organizers}

The derived category of projective space is among the best understood objects in
algebraic geometry. Beilinson's resolution of the diagonal shows that the line bundles
$\cO(-n),\dots,\cO$ form a full strong exceptional collection for $\Db(\PP^n)$, so that
every object of $\Db(\PP^n)$ is represented, uniquely up to homotopy, by a minimal complex
whose terms are sums of these line bundles~\cite{Bei78}. In a complementary direction, the
Bernstein--Gel'fand--Gel'fand correspondence~\cite{BGG78} and the Tate resolutions of
Eisenbud, Fl{\o}ystad, and Schreyer~\cite{EFS03} identify $\Db(\PP^n)$ with a homotopy
category of exact complexes over an exterior algebra, in such a way that the terms of the
complex attached to a sheaf record the cohomology of all of its twists.

Toric varieties are the natural next testing ground for these ideas. By work of
Cox~\cite{Cox95}, a toric variety $X$ is a quotient of an open subset of affine space by a
group determined by the $\Cl(X)$-grading on its Cox ring $S$, and coherent sheaves on $X$
(or on the associated toric stack, when $X$ is only simplicial) correspond to
$\Cl(X)$-graded $S$-modules modulo torsion with respect to the irrelevant ideal. Analogues
of Beilinson's theorem nevertheless proved elusive. King conjectured in 1997 that every
smooth projective toric variety admits a full strong exceptional collection of line
bundles~\cite{Kin97}; this was disproved by Hille and Perling~\cite{HP06}, and further
counterexamples of Micha{\l}ek and Efimov~\cite{Mic11,Efi14} showed that the conjecture
cannot be rescued by natural restrictions such as the Fano condition. The closest positive
result, due to Kawamata~\cite{Kaw06}, produces full exceptional collections that are
neither strong nor made of line bundles.

The last few years have seen rapid progress, and the aim of this Arbeitsgemeinschaft was
to give a complete account of three interlocking developments.

\subsection*{The Hanlon--Hicks--Lazarev resolution}
Hanlon, Hicks, and Lazarev~\cite{HHL24} constructed, for any toric subvariety of a toric
variety, an explicit resolution by line bundles. In particular, the diagonal of a smooth
toric variety $X$ has a resolution of length $\dim X$ whose terms are sums of line bundles
$L\boxtimes L'$ with $L$ and $L'$ in the Bondal--Thomsen collection, i.e.\ among the line
bundles appearing as summands of toric Frobenius pushforwards of
$\cO_X$~\cite{Bon06,Tho00}. The resolution is cellular: it is supported on the subdivision
of a real torus by coordinate hyperplanes, with labels obtained by taking ceilings of
coordinates. The original proof was symplectic in origin; recent work of Berkesch,
Cranton Heller, Smith, and Yang~\cite{BCHSY25} and of Borisov and Han~\cite{BH25} gave
short algebraic proofs, interpreting the construction as a resolution of the
normalization of a toric ring. These resolutions quickly yielded proofs of three
conjectures: Bondal's conjecture that the Bondal--Thomsen collection generates
$\Db(X)$~\cite{Bon06}; the conjecture of Berkesch, Erman, and Smith that every finitely
generated module over the Cox ring admits a virtual resolution of length at most
$\dim X$~\cite{BES20}; and the toric case of Orlov's conjecture that the Rouquier
dimension of $\Db(X)$ equals $\dim X$~\cite{Orl09}, proven independently by Favero and
Huang~\cite{FH23}.

\subsection*{The Cox category and a revised King's conjecture}
A striking feature of the Bondal--Thomsen collection and of the Hanlon--Hicks--Lazarev
resolution of the diagonal is that both depend only on the graded Cox ring $S$, and not on
the fan. The toric varieties (and stacks) sharing a given Cox ring are parametrized by the
chambers of its secondary fan, which suggests that the right analogue of Beilinson's
theorem should see all of them at once. In joint work of Ballard, Berkesch, Brown,
Cranton Heller, Erman, Favero, Ganatra, Hanlon, and Huang~\cite{BBBplus24}, this is made
precise through the \emph{Cox category} $\DCox(S)$: the subcategory of the derived category
of a common toric stack resolution generated by the pullbacks of the derived categories of
all the toric stacks $\cX_1,\dots,\cX_r$ arising from the secondary fan. The main
theorem, a revised King's conjecture, states that the Bondal--Thomsen collection, with
each line bundle pulled back from the chamber containing its degree, is a full strong
exceptional collection for $\DCox(S)$. The key technical input is a
\emph{$\Theta$-transform lemma}: Fourier--Mukai transforms between chambers send
Bondal--Thomsen line bundles to the corresponding Bondal--Thomsen line bundles, with no
higher direct images. Consequences include that $\DCox(S)$ is smooth and proper of
Rouquier dimension $\dim X$; an equivalence between $\DCox(S)$ and an explicit homotopy
category of complexes of free $S$-modules; a resolution of the diagonal for $\DCox(S)$
explaining the uniformity of the Hanlon--Hicks--Lazarev resolution across the secondary
fan; and noncommutative resolutions of singularities for semiprojective toric varieties,
extending the affine results of \v{S}penko--Van den Bergh~\cite{SVdB17} and
Faber--Muller--Smith~\cite{FMS19}.

\subsection*{Tate resolutions on toric varieties}
Building on the case of products of projective spaces~\cite{EES15}, Brown and Erman
extended the BGG correspondence and Tate resolutions to toric varieties~\cite{BE26}.
Because the $\Cl(X)$-grading no longer induces a homological grading, the exterior side
must be formulated in terms of differential modules. The toric Tate resolution of a
coherent sheaf $\cF$ is a minimal free differential module over the exterior algebra whose Betti numbers 
encode the cohomology table of $\cF$. This leads to a toric BGG equivalence between $\Db(\cX)$ and a
homotopy category of free differential modules~\cite{BEkoszul}, as well as to
an algorithm for computing sheaf cohomology on weighted projective stacks~\cite{BE25}.

\subsection*{Course of the workshop}
Following the Arbeitsgemeinschaft format, all nineteen talks were given by participants,
following a detailed program and set of lecture notes prepared by the organizers together
with Lauren Cranton Heller, Mahrud Sayrafi, and Jay Yang. Several running examples
recurred throughout the week, chief among them the Hirzebruch surface $\mathbb{H}_3$,
whose secondary fan has two chambers, corresponding to $\mathbb{H}_3$ and the weighted
projective stack $\PP(1,1,3)$; it is the smallest example in which the Cox category is a
genuine combination of two derived categories.

Monday was devoted to background. Three talks treated $\Db(\PP^n)$: derived categories
and Fourier--Mukai transforms; Beilinson's resolution of the diagonal and full strong
exceptional collections; and the BGG correspondence, culminating in the identification of
$\Db(\PP^n)$ with the singularity category of the exterior algebra. Two further talks
introduced toric varieties, Cox rings, the toric module--sheaf correspondence, and toric stacks.

On Tuesday, a talk on the Bondal--Thomsen collection (via the Bondal zonotope, Frobenius
pushforwards, and Musta\c{t}\u{a} vanishing) was followed by three talks on the
Hanlon--Hicks--Lazarev resolution: its construction and first examples, further features
including the resolution of the diagonal of $\mathbb{H}_3$, and a complete algebraic
proof. Several talks emphasized computation of these resolutions using Macaulay2. 

Wednesday morning covered the applications of the Hanlon--Hicks--Lazarev theorem to the
conjectures of Bondal, Berkesch--Erman--Smith, and Orlov; the secondary fan, from the
points of view of GIT, polytopes, and irrelevant ideals; and the definition of the Cox
category. The afternoon was reserved for the traditional excursion.

On Thursday, a talk on Fourier--Mukai transforms between the toric stacks of the secondary
fan was followed by a three-part proof of the revised King's conjecture: an outline
reducing orthogonality to the $\Theta$-transform lemma; an analysis of that lemma,
including the Atiyah flop as an illustration of the delicate balance between the
$R^0$-term and the vanishing of higher direct images; and the proof of fullness. %The day
closed with an open question-and-answer session.

Friday featured consequences of the revised King's conjecture (smoothness, properness,
Rouquier dimension, and noncommutative resolutions), followed by two talks on Tate
resolutions: first over $\PP^n$, and then over toric varieties via differential modules.

% TODO (organizers): optionally add a sentence on participation and on any
% notable outcomes of the discussion sessions.

Several open directions emerged naturally from the program: whether the Cox category, or
invariants such as its K-theory and Hochschild homology, admit an intrinsic geometric
description; whether the $\Theta$-transform lemma admits a simpler proof, for instance
via vanishing theorems from birational geometry; a geometric description of the
differentials in Tate resolutions; and whether toric Tate resolutions can be constructed
from ordinary free resolutions when the class group has rank greater than one.

\newpage
\bibliographystyle{amsalpha}
\bibliography{owr_introduction}

\end{document}
